\scanpage{372}\mspage{288}
sous-catégories \unsure{strictemt|strictmt|strictement} pleines \inserted{$S(\pi)$} de $M_\pi$
qui sont \deleted{\unsure{stables|sta}}\unsure{karoubiennes}, ou ce qui
revient au $=$, stables \deleted{par} dans $M_\pi$ par
facteurs directs. Soit $\Sigma'_\pi$ l'ens
des sous-\unsure{catégorie|catégories} de $M$ qui sont
$\pi$-accessibles dans $M$. Considérons
les applications
\diagram{\begin{tikzcd}[ampersand replacement=\&]
\text{(4.13.4.1)} \& \Sigma_\pi \arrow[r,shift left=.7ex,"\alpha_\pi"] \& {\Sigma'_\pi} \arrow[l,shift left=.7ex,"\beta_\pi"]
\end{tikzcd}}
définies ainsi : Si $S(\pi)\deleted{\text{\unsure{est}}}\in\Sigma_\pi$,
$\alpha_\pi(S(\pi))$ est l'image essentielle du
foncteur \inserted{\unsure{can} pl fid} $\mathrm{Ind}_\pi(S(\pi))\unsure{\longrightarrow|\text{---}} M$.
Si $\underline S\in\Sigma'_\pi$, $\beta_\pi(\underline S)=\underline S\cap M_\pi$.
Les applications $\alpha_\pi,\beta_\pi$ sont des
bijections, inverses l'une de
l'autre. \deleted{De plus, si $\underline S$ et $S$}

\deleted{\unsure{On n'a pas}} \inserted{On n'a pas} eu à supposer que
$\pi$ soit utile pour $M$, ni que \deleted{\unsure{l'on}}$\pi$\deleted{\unsure{ait}}
\deleted{ait $\pi\geq\pi_0$, où $\pi_0$ est le plus petit
cardinal utile} \inserted{soit \unsure{minoré|majoré} par un cardinal utile pour $M$,}
i.e. qu'on ait $\pi\geq\pi_0$, où
$\pi_0$ est le plus petit cardinal utile pour $M$
% archive sheet number: 371
